% Reusable table and figure templates extracted from the AAIB research report.
% Copy this file into a new LaTeX research repo and adapt the placeholders.
% Requires: float, array, tabularx, longtable, booktabs, tikz, pgfplots, xcolor, patterns.

% ---------------------------------------------------------------------------
% Required palette and package assumptions
% ---------------------------------------------------------------------------
% \usepackage{float}
% \usepackage{array}
% \usepackage{tabularx}
% \usepackage{longtable}
% \usepackage{booktabs}
% \usepackage{tikz}
% \usetikzlibrary{patterns}
% \usepackage{pgfplots}
% \pgfplotsset{compat=1.18}
% \usepackage{xcolor}
%
% \definecolor{aaibblue}{RGB}{0,71,131}
% \definecolor{aaibgold}{RGB}{185,150,60}
% \definecolor{aaibgray}{RGB}{120,120,120}
% \definecolor{aaibred}{RGB}{180,40,40}
% \definecolor{aaibgreen}{RGB}{34,120,60}
% \definecolor{aaiblightblue}{RGB}{100,160,210}

% ---------------------------------------------------------------------------
% Table template 1: three-column stakeholder / summary table
% ---------------------------------------------------------------------------
\begin{table}[H]
\centering
\caption{<Descriptive Table Caption>}
\label{tab:<short-label>}
\small
\begin{tabularx}{\textwidth}{>{\raggedright\arraybackslash}X
                >{\raggedright\arraybackslash}X
                >{\raggedright\arraybackslash}X}
\toprule
\textbf{<Column 1>} & \textbf{<Column 2>} & \textbf{<Column 3>} \\
\midrule
<Row label> & <Short description> & <Primary interpretation or implication> \\
<Row label> & <Short description> & <Primary interpretation or implication> \\
<Row label> & <Short description> & <Primary interpretation or implication> \\
\bottomrule
\end{tabularx}
\end{table}

% ---------------------------------------------------------------------------
% Table template 2: four-column risk / criteria / monitoring table
% ---------------------------------------------------------------------------
\begin{table}[H]
\centering
\caption{<Descriptive Table Caption>}
\label{tab:<short-label>}
\small
\begin{tabularx}{\textwidth}{>{\raggedright\arraybackslash}X
                l
                >{\raggedright\arraybackslash}X
                >{\raggedright\arraybackslash}X}
\toprule
\textbf{<Factor>} & \textbf{<Rating>} & \textbf{<Evidence or Threshold>} & \textbf{<Mitigant or Implication>} \\
\midrule
<Factor name> & <High/Medium/Low> & <Metric, basis, or threshold> & <Implication> \\
<Factor name> & <High/Medium/Low> & <Metric, basis, or threshold> & <Implication> \\
<Factor name> & <High/Medium/Low> & <Metric, basis, or threshold> & <Implication> \\
\bottomrule
\end{tabularx}
\end{table}

% ---------------------------------------------------------------------------
% Table template 3: ranked market / peer comparison table
% ---------------------------------------------------------------------------
\begin{table}[H]
\centering
\caption{<Descriptive Table Caption>}
\label{tab:<short-label>}
\small
\begin{tabularx}{\textwidth}{c >{\raggedright\arraybackslash}X r >{\raggedright\arraybackslash}X}
\toprule
\textbf{Rank} & \textbf{Entity} & \textbf{Metric} & \textbf{Type or Note} \\
\midrule
1 & <Entity name> & <Value> & <Classification or note> \\
2 & <Entity name> & <Value> & <Classification or note> \\
3 & \textbf{<Focal entity>} & \textbf{<Value>} & \textbf{<Classification or note>} \\
\bottomrule
\end{tabularx}
\end{table}

% Source: <Source name, date>.

% ---------------------------------------------------------------------------
% Table template 4: compact numeric comparison table
% ---------------------------------------------------------------------------
\begin{table}[H]
\centering
\caption{<Descriptive Table Caption>}
\label{tab:<short-label>}
\small
\begin{tabular}{lrrr}
\toprule
\textbf{Metric} & \textbf{<Period 1>} & \textbf{<Period 2>} & \textbf{Change} \\
\midrule
<Metric name> & <Value> & <Value> & <Change> \\
<Metric name> & <Value> & <Value> & <Change> \\
<Metric name> & <Value> & <Value> & <Change> \\
\bottomrule
\end{tabular}
\end{table}

% ---------------------------------------------------------------------------
% Table template 5: wide digital / capability comparison table
% ---------------------------------------------------------------------------
\begin{table}[H]
\centering
\caption{<Descriptive Table Caption>}
\label{tab:<short-label>}
\small
\begin{tabularx}{\textwidth}{l >{\raggedright\arraybackslash}X
                >{\raggedright\arraybackslash}X
                >{\raggedright\arraybackslash}X}
\toprule
\textbf{Dimension} & \textbf{<Entity A>} & \textbf{<Entity B>} & \textbf{<Entity C>} \\
\midrule
<Dimension> & <Assessment or metric> & <Assessment or metric> & <Assessment or metric> \\
<Dimension> & <Assessment or metric> & <Assessment or metric> & <Assessment or metric> \\
<Dimension> & <Assessment or metric> & <Assessment or metric> & <Assessment or metric> \\
\bottomrule
\end{tabularx}
\end{table}

% Source: <Source list>.

% ---------------------------------------------------------------------------
% Table template 6: five-column landscape / ecosystem table
% ---------------------------------------------------------------------------
\begin{table}[H]
\centering
\caption{<Descriptive Table Caption>}
\label{tab:<short-label>}
\small
\begin{tabularx}{\textwidth}{l >{\raggedright\arraybackslash}X l l >{\raggedright\arraybackslash}X}
\toprule
\textbf{Entity} & \textbf{Model} & \textbf{Funding or Size} & \textbf{Scale} & \textbf{Strategic Overlap} \\
\midrule
<Entity> & <Business model> & <Value> & <Users/assets/reach> & <Overlap or implication> \\
<Entity> & <Business model> & <Value> & <Users/assets/reach> & <Overlap or implication> \\
<Entity> & <Business model> & <Value> & <Users/assets/reach> & <Overlap or implication> \\
\bottomrule
\end{tabularx}
\end{table}

% ---------------------------------------------------------------------------
% Longtable template 1: appendix financial highlights
% ---------------------------------------------------------------------------
{
\scriptsize
\setlength{\tabcolsep}{2pt}
\begin{longtable}[]{@{}
  >{\raggedright\arraybackslash}p{0.26\linewidth}
  >{\raggedright\arraybackslash}p{0.11\linewidth}
  >{\raggedright\arraybackslash}p{0.11\linewidth}
  >{\raggedright\arraybackslash}p{0.11\linewidth}
  >{\raggedright\arraybackslash}p{0.11\linewidth}@{}}
\caption{<Appendix Table Caption>} \label{tab:<short-label>} \\
\toprule\noalign{}
Metric & <Year 1> & <Year 2> & <Year 3> & <Year 4> \\
\midrule\noalign{}
\endfirsthead
\caption[]{<Appendix Table Caption>, continued} \\
\toprule\noalign{}
Metric & <Year 1> & <Year 2> & <Year 3> & <Year 4> \\
\midrule\noalign{}
\endhead
\bottomrule\noalign{}
\endlastfoot
<Metric name> & <Value> & <Value> & <Value> & <Value> \\
<Metric name> & <Value> & <Value> & <Value> & <Value> \\
<Metric name> & <Value> & <Value> & <Value> & <Value> \\
\end{longtable}
\normalsize
}

% ---------------------------------------------------------------------------
% Longtable template 2: appendix timeline table
% ---------------------------------------------------------------------------
{
\footnotesize
\renewcommand{\arraystretch}{1.15}
\begin{longtable}[]{@{}
  >{\raggedright\arraybackslash}p{0.15\linewidth}
  >{\raggedright\arraybackslash}p{0.79\linewidth}@{}}
\caption{<Timeline Caption>} \label{tab:<short-label>} \\
\toprule\noalign{}
Year & Milestone \\
\midrule\noalign{}
\endfirsthead
\caption[]{<Timeline Caption>, continued} \\
\toprule\noalign{}
Year & Milestone \\
\midrule\noalign{}
\endhead
\bottomrule\noalign{}
\endlastfoot
<Year> & <Milestone description> \\
<Year> & <Milestone description> \\
<Year> & <Milestone description> \\
\end{longtable}
\normalsize
}

% ---------------------------------------------------------------------------
% Figure template 1: line chart with shaded area
% ---------------------------------------------------------------------------
\begin{figure}[H]
  \centering
  \begin{tikzpicture}
    \begin{axis}[
      width=0.85\textwidth,
      height=6cm,
      ylabel={<Y-axis label>},
      xmin=<xmin>, xmax=<xmax>,
      ymin=0, ymax=<ymax>,
      xtick={<x ticks>},
      ytick={<y ticks>},
      ymajorgrids=true,
      grid style={gray!20, dashed},
      axis x line=bottom,
      axis y line=left,
      tick align=outside,
      tick label style={font=\footnotesize},
      label style={font=\small},
      clip=false,
    ]
    \addplot[
      fill=aaibblue!15,
      draw=none,
    ] coordinates {
      (<x1>, <y1>)
      (<x2>, <y2>)
      (<x3>, <y3>)
      (<x3>, 0)
      (<x1>, 0)
    } \closedcycle;

    \addplot[
      color=aaibblue,
      line width=1.8pt,
      mark=square*,
      mark size=3pt,
      mark options={fill=aaibblue},
    ] coordinates {
      (<x1>, <y1>)
      (<x2>, <y2>)
      (<x3>, <y3>)
    };

    \node[font=\footnotesize, text=aaibblue, above left] at (axis cs:<x1>,<y1>) {<label>};
    \node[font=\footnotesize, text=aaibblue, above right] at (axis cs:<x3>,<y3>) {<label>};
    \end{axis}
  \end{tikzpicture}
  \vspace{4pt}
  \caption{<Chart title, period (unit). Analytical takeaway sentence. Source: source name.>}
  \label{fig:<short-label>}
\end{figure}

% ---------------------------------------------------------------------------
% Figure template 2: vertical bar chart
% ---------------------------------------------------------------------------
\begin{figure}[H]
  \centering
  \begin{tikzpicture}
    \begin{axis}[
      width=0.78\textwidth,
      height=6.5cm,
      ybar,
      bar width=22pt,
      ylabel={<Y-axis label>},
      symbolic x coords={<A>,<B>,<C>,<D>},
      xtick=data,
      ymin=0, ymax=<ymax>,
      ytick={<y ticks>},
      ymajorgrids=true,
      grid style={gray!20, dashed},
      axis x line=bottom,
      axis y line=left,
      tick align=outside,
      tick label style={font=\footnotesize},
      label style={font=\small},
      nodes near coords,
      nodes near coords style={font=\footnotesize, color=aaibblue, anchor=south},
      every node near coord/.append style={yshift=2pt},
      enlarge x limits=0.2,
      clip=false,
    ]
    \addplot[fill=aaibblue, draw=aaibblue!80!black] coordinates {
      (<A>, <value>)
      (<B>, <value>)
      (<C>, <value>)
      (<D>, <value>)
    };
    \end{axis}
  \end{tikzpicture}
  \vspace{4pt}
  \caption{<Chart title, period (unit). Analytical takeaway sentence. Source: source name.>}
  \label{fig:<short-label>}
\end{figure}

% ---------------------------------------------------------------------------
% Figure template 3: grouped vertical bar chart, two periods
% ---------------------------------------------------------------------------
\begin{figure}[H]
  \centering
  \begin{tikzpicture}
    \begin{axis}[
      width=0.82\textwidth,
      height=6.5cm,
      ybar=4pt,
      bar width=16pt,
      ylabel={<Y-axis label>},
      symbolic x coords={<Category 1>,<Category 2>,<Category 3>},
      xtick=data,
      xticklabel style={font=\footnotesize},
      ymin=0, ymax=<ymax>,
      ytick={<y ticks>},
      ymajorgrids=true,
      grid style={gray!20, dashed},
      axis x line=bottom,
      axis y line=left,
      tick align=outside,
      tick label style={font=\footnotesize},
      label style={font=\small},
      legend style={at={(0.5,-0.18)},anchor=north,legend columns=2,font=\footnotesize},
      nodes near coords,
      nodes near coords style={font=\footnotesize, anchor=south},
      every node near coord/.append style={yshift=2pt},
      enlarge x limits=0.25,
      clip=false,
    ]
    \addplot[fill=aaiblightblue!60, draw=aaiblightblue!80!black] coordinates {
      (<Category 1>, <value>)
      (<Category 2>, <value>)
      (<Category 3>, <value>)
    };
    \addlegendentry{<Period 1>}

    \addplot[fill=aaibblue, draw=aaibblue!80!black] coordinates {
      (<Category 1>, <value>)
      (<Category 2>, <value>)
      (<Category 3>, <value>)
    };
    \addlegendentry{<Period 2>}
    \end{axis}
  \end{tikzpicture}
  \vspace{4pt}
  \caption{<Chart title, periods (unit). Analytical takeaway sentence. Source: source name.>}
  \label{fig:<short-label>}
\end{figure}

% ---------------------------------------------------------------------------
% Figure template 4: line chart with threshold
% ---------------------------------------------------------------------------
\begin{figure}[H]
  \centering
  \begin{tikzpicture}
    \begin{axis}[
      width=0.85\textwidth,
      height=6cm,
      ylabel={<Y-axis label>},
      xmin=<xmin>, xmax=<xmax>,
      ymin=0, ymax=<ymax>,
      xtick={<x ticks>},
      ytick={<y ticks>},
      ymajorgrids=true,
      grid style={gray!20, dashed},
      axis x line=bottom,
      axis y line=left,
      tick align=outside,
      tick label style={font=\footnotesize},
      label style={font=\small},
      legend style={at={(0.97,0.97)},anchor=north east,legend columns=1,font=\footnotesize,fill=white,draw=gray!40},
      clip=false,
    ]
    \addplot[
      color=aaibblue,
      line width=1.8pt,
      mark=*,
      mark size=3pt,
      mark options={fill=aaibblue},
    ] coordinates {
      (<x1>, <y1>)
      (<x2>, <y2>)
      (<x3>, <y3>)
    };
    \addlegendentry{<Main metric>}

    \addplot[
      color=aaibred,
      line width=1pt,
      dashed,
      domain=<xmin>:<xmax>,
    ] {<threshold>};
    \addlegendentry{<Threshold label>}

    \node[font=\footnotesize, text=aaibblue, above right] at (axis cs:<x1>,<y1>) {<label>};
    \node[font=\footnotesize, text=aaibblue, above right] at (axis cs:<x3>,<y3>) {<label>};
    \end{axis}
  \end{tikzpicture}
  \vspace{4pt}
  \caption{<Chart title, period. Analytical takeaway sentence. Source: source name.>}
  \label{fig:<short-label>}
\end{figure}

% ---------------------------------------------------------------------------
% Figure template 5: horizontal peer-ranking bar chart with focal highlight
% ---------------------------------------------------------------------------
\begin{figure}[H]
  \centering
  \begin{tikzpicture}
    \begin{axis}[
      width=0.85\textwidth,
      height=5.8cm,
      xbar=0pt,
      bar width=15pt,
      xlabel={<X-axis label>},
      ylabel={},
      symbolic y coords={<Focal>,<Peer 1>,<Peer 2>,<Peer 3>,<Peer 4>},
      ytick={<Focal>,<Peer 1>,<Peer 2>,<Peer 3>,<Peer 4>},
      xmin=<xmin>, xmax=<xmax>,
      xtick={<x ticks>},
      xmajorgrids=true,
      grid style={gray!20, dashed},
      axis x line=bottom,
      axis y line=left,
      tick align=outside,
      tick label style={font=\footnotesize},
      y tick label style={font=\footnotesize},
      label style={font=\small},
      nodes near coords,
      nodes near coords style={font=\footnotesize, anchor=west},
      nodes near coords align={horizontal},
      point meta=explicit symbolic,
      enlarge y limits=0.18,
      clip=false,
    ]
    \addplot[fill=aaiblightblue!55, draw=aaiblightblue!80!black, bar shift=0pt] coordinates {
      (<value>,<Peer 1>) [{<label>}]
      (<value>,<Peer 2>) [{<label>}]
      (<value>,<Peer 3>) [{<label>}]
      (<value>,<Peer 4>) [{<label>}]
    };

    \addplot[fill=aaibblue, draw=aaibblue!80!black, bar shift=0pt] coordinates {
      (<value>,<Focal>) [{<label>}]
    };
    \end{axis}
  \end{tikzpicture}
  \vspace{4pt}
  \caption{<Chart title, period (unit). Analytical takeaway sentence. Source: source name.>}
  \label{fig:<short-label>}
\end{figure}

% ---------------------------------------------------------------------------
% Figure template 6: horizontal actual-vs-minimum bar chart
% ---------------------------------------------------------------------------
\begin{figure}[H]
  \centering
  \begin{tikzpicture}
    \begin{axis}[
      width=0.85\textwidth,
      height=5.5cm,
      xbar=3pt,
      bar width=13pt,
      xlabel={<X-axis label>},
      ylabel={},
      symbolic y coords={<Metric 3>,<Metric 2>,<Metric 1>},
      ytick=data,
      y tick label style={font=\footnotesize},
      xmin=0, xmax=<xmax>,
      xtick={<x ticks>},
      xmajorgrids=true,
      grid style={gray!20, dashed},
      axis x line=bottom,
      axis y line=left,
      tick align=outside,
      tick label style={font=\footnotesize},
      label style={font=\small},
      legend style={at={(0.5,-0.32)},anchor=north,legend columns=2,font=\footnotesize},
      nodes near coords,
      nodes near coords style={font=\footnotesize, anchor=west},
      nodes near coords align={horizontal},
      point meta=explicit symbolic,
      enlarge y limits=0.25,
      clip=false,
    ]
    \addplot[
      fill=aaibred!30,
      draw=aaibred,
      pattern=north east lines,
      pattern color=aaibred,
    ] coordinates {
      (<minimum>,<Metric 1>) [{<label>}]
      (<minimum>,<Metric 2>) [{<label>}]
      (<minimum>,<Metric 3>) [{<label>}]
    };
    \addlegendentry{<Minimum or threshold>}

    \addplot[fill=aaibblue, draw=aaibblue!80!black] coordinates {
      (<actual>,<Metric 1>) [{<label>}]
      (<actual>,<Metric 2>) [{<label>}]
      (<actual>,<Metric 3>) [{<label>}]
    };
    \addlegendentry{<Actual series>}
    \end{axis}
  \end{tikzpicture}
  \vspace{4pt}
  \caption{<Chart title, period (unit). Analytical takeaway sentence. Source: source name.>}
  \label{fig:<short-label>}
\end{figure}

% ---------------------------------------------------------------------------
% Figure template 7: vertical bar chart with category allocation
% ---------------------------------------------------------------------------
\begin{figure}[H]
  \centering
  \begin{tikzpicture}
    \begin{axis}[
      width=0.85\textwidth,
      height=6.2cm,
      ybar,
      bar width=18pt,
      ylabel={<Y-axis label>},
      symbolic x coords={<Category 1>,<Category 2>,<Category 3>,<Category 4>},
      xtick=data,
      xticklabel style={font=\footnotesize, text width=2.2cm, align=center},
      ymin=0, ymax=<ymax>,
      ytick={<y ticks>},
      ymajorgrids=true,
      grid style={gray!20, dashed},
      axis x line=bottom,
      axis y line=left,
      tick align=outside,
      tick label style={font=\footnotesize},
      label style={font=\small},
      nodes near coords,
      nodes near coords style={font=\footnotesize, color=aaibblue, anchor=south},
      every node near coord/.append style={yshift=2pt},
      enlarge x limits=0.18,
      clip=false,
    ]
    \addplot[
      fill=aaibblue,
      draw=aaibblue!80!black,
    ] coordinates {
      (<Category 1>, <value>)
      (<Category 2>, <value>)
      (<Category 3>, <value>)
      (<Category 4>, <value>)
    };
    \end{axis}
  \end{tikzpicture}
  \vspace{4pt}
  \caption{<Chart title, period (unit). Analytical takeaway sentence. Source: source name.>}
  \label{fig:<short-label>}
\end{figure}
